% =============================================================================
% AI and Machine Learning -- Project Report Template
% Alpen-Adria-Universität Klagenfurt
%
% Page limit: 9 pages, EXCLUDING references. Everything else counts,
% including figures, tables, the appendix and the AI disclosure statement.
% Do not change margins, font sizes or spacing.
%
% The comments under each heading describe what the section must contain.
% Replace every "Your text." with your content and submit the PDF in Moodle,
% together with your code (see the project guide, README.md).
%
% State in every section who did what, e.g. "The logistic regression baseline
% was implemented and evaluated by (name) ...".
% =============================================================================

\documentclass{article}

% "preprint" shows your names; without it the template produces an
% anonymous submission.
\usepackage[preprint]{neurips_2025}

\usepackage[utf8]{inputenc} % allow utf-8 input
\usepackage[T1]{fontenc}    % use 8-bit T1 fonts
\usepackage{hyperref}       % hyperlinks
\usepackage{url}            % simple URL typesetting
\usepackage{booktabs}       % professional-quality tables
\usepackage{amsmath}        % mathematical environments
\usepackage{amsfonts}       % blackboard math symbols
\usepackage{nicefrac}       % compact symbols for 1/2, etc.
\usepackage{microtype}      % microtypography
\usepackage{xcolor}         % colors
\usepackage{graphicx}       % figures

\title{Your Project Title}

\author{%
  Student A \\
  Student ID: 00000000 \\
  \texttt{student.a@edu.aau.at} \\
  \And
  Student B \\
  Student ID: 00000000 \\
  \texttt{student.b@edu.aau.at} \\
  \And
  Student C \\
  Student ID: 00000000 \\
  \texttt{student.c@edu.aau.at} \\
}

\begin{document}

\maketitle

\begin{abstract}
% Brief overview of the domain problem, your methodology and the key empirical
% findings.
Your abstract.
\end{abstract}

% =============================================================================
\section{Motivation}
% - The domain problem, why it matters, and your hypothesis.
% - A short summary of your approach and results.

Your text.

% =============================================================================
\section{Data Exploration and Preprocessing}
% - Exploratory data analysis: distributions, correlations, missing values and
%   class imbalance.
% - Feature engineering, encoding and scaling, encapsulated in leak-free
%   sklearn.pipeline.Pipeline objects.

Your text.

% =============================================================================
\section{Theoretical Background and Model Selection}
% - Mathematical formulation of each chosen model family (one per student).
% - Your hypotheses on why each algorithm suits the characteristics of the data.

Your text.

% =============================================================================
\section{Experimental Setup and Hyperparameter Optimization}
% - Validation methodology: stratified K-fold cross-validation or a temporal
%   split.
% - Hyperparameter search spaces and optimization strategies (grid search,
%   random search, Bayesian optimization).
% - Links to your code and to your public W&B project (see the paragraph
%   below).

Your text.

\paragraph{Code and experiments.}
Code: \url{https://github.com/your-team/your-project} (or a Colab notebook). \\
Experiment tracking: \url{https://wandb.ai/your-entity/your-project}.

% =============================================================================
\section{Evaluation and Diagnostics}
% - One consolidated table comparing the models of all team members with the
%   literature baselines (see Table~\ref{tab:results}).
% - Diagnostic figures: ROC curves, precision-recall curves, confusion matrices
%   and feature importance rankings (see Figure~\ref{fig:diagnostics}).
% - Results and figures should come from runs logged in your public W&B project.

Your text.

\begin{table}[t]
  \caption{Test results of all models and the literature baselines (mean $\pm$ standard deviation over the cross-validation folds).}
  \label{tab:results}
  \centering
  \begin{tabular}{llccc}
    \toprule
    Model & Team member & $F_1$-score & ROC-AUC & Runtime \\
    \midrule
    Literature baseline (cite source)       & ---       & 0.00 & 0.00 & --- \\
    \midrule
    e.g.\ logistic regression               & Student A & 0.00 $\pm$ 0.00 & 0.00 $\pm$ 0.00 & 0\,s \\
    e.g.\ gradient boosting                 & Student B & 0.00 $\pm$ 0.00 & 0.00 $\pm$ 0.00 & 0\,s \\
    e.g.\ SVM (RBF kernel)                  & Student C & 0.00 $\pm$ 0.00 & 0.00 $\pm$ 0.00 & 0\,s \\
    \bottomrule
  \end{tabular}
\end{table}

\begin{figure}[t]
  \centering
  % Replace the box with your figure, e.g.:
  % \includegraphics[width=0.8\linewidth]{figures/roc_curves.pdf}
  \fbox{\parbox[c][4cm][c]{0.8\linewidth}{\centering ROC and precision-recall curves of all models}}
  \caption{Diagnostic curves of all models.}
  \label{fig:diagnostics}
\end{figure}

% =============================================================================
\section{Error Analysis and Discussion}
% - Failure modes: which instances did the models misclassify, and why?
% - Domain interpretation: which actionable recommendations follow from the
%   models?

Your text.

% =============================================================================
\section{Conclusions and Future Directions}
% - Summary of your findings.
% - Recommendations for deployment or further research.

Your text.

% =============================================================================
\bibliographystyle{plainnat}
\bibliography{references}

% =============================================================================
\appendix

\section{AI Disclosure Statement}
% - Declare any use of AI tools, e.g. for brainstorming, exploratory data
%   analysis, code generation, text drafting or literature search.
% - Specify the tools, model versions and the general purpose of their use.
% - If you did not use any, say so.

Your text.

\end{document}
